\chapter{Routing Replay and Mechanism Studies}
\label{chap:routing-studies}
\chapterrunning{routing studies}

The routing phases move the project from behavioral game environments to benchmark optimization.
This transition matters because routing benchmarks provide stronger baselines, fixed held-out
panels, and more standardized notions of progress. The original routing hypothesis was that
replay-aware code evolution would improve held-out benchmark gap, and that the improvement would be
strongest when replay focused on informative failures. The completed evidence is more nuanced.

\section{Phase-5 Routing Transfer Across TSP, ATSP, and CVRP}
\label{sec:routing-phase5}

The first routing campaign compares four replay conditions across three benchmark families:

\begin{enumerate}[leftmargin=2em]
\item no replay,
\item random replay,
\item raw-gap failure replay, and
\item failure replay with compression pressure.
\end{enumerate}

Each family uses twenty paired offsets and reports paired bootstrap intervals on the relevant held
out benchmark endpoint. The resulting evidence is summarized in
Table~\ref{tab:routing-phase5-summary}.

\begin{table}[t]
\centering
\scriptsize
\setlength{\tabcolsep}{2pt}
\begin{tabularx}{\textwidth}{
  >{\raggedright\arraybackslash}p{0.58in}
  >{\raggedright\arraybackslash}p{0.9in}
  Y
  >{\raggedright\arraybackslash}p{0.95in}
  Y
}
\toprule
Family & Best mean condition & Key paired comparison & 95\% interval & Interpretation \\
\midrule
TSP & random replay & held-out TSPLIB delta versus no replay: $-0.0311$; combined transfer delta: $-0.0282$ & $[-0.0602,-0.0008]$ on TSPLIB; $[-0.0520,-0.0029]$ on combined transfer & clear positive replay signal, but the winning arm is random rather than failure replay \\
ATSP & failure replay with compression & held-out ATSP delta versus no replay: $-0.0056$; combined transfer delta: $-0.0063$ & $[-0.0191,0.0071]$ on ATSP; $[-0.0133,-0.0002]$ on combined transfer & weakly favorable, not decisively separated on the primary endpoint \\
CVRP & no replay & failure replay versus random replay on held-out CVRPLIB: $-0.0076$ & $[-0.0146,-0.0006]$ & replay is not dominant; no-replay remains the strongest family-level baseline \\
\bottomrule
\end{tabularx}
\caption{Phase-5 family summaries across the first replicated routing campaign. Lower deltas are
better because all routing endpoints are lower-is-better gap metrics.}
\label{tab:routing-phase5-summary}
\end{table}

Table~\ref{tab:routing-phase5-summary} rules out a simple replay-dominance claim. Replay does help
on symmetric TSP, but the strongest arm is random replay rather than raw failure replay. ATSP
retains only a weak compression-aware signal, and the primary held-out ATSP endpoint is not cleanly
separated. In the replicated phase-5 CVRP suite, no-replay remains the best overall condition.

The practical consequence is important. If replay had behaved uniformly across TSP, ATSP, and
CVRP, then the thesis could argue for a reusable memory mechanism with broad routing scope. The
actual outcome is narrower: replay pressure interacts with benchmark structure, and the initially
favored raw-failure recipe is not the stable winner even within routing.

\section{Phase-6 TSP Replay Mechanism Study}
\label{sec:routing-phase6}

Phase 6 asks why random replay beat raw failure replay in the initial TSP study. The follow-up
campaign remains within symmetric TSPLIB TSP, but expands the mechanism space to curated and
diagnostic replay variants, including diversity-aware failure replay, residual-failure replay, and
compression applied to the stronger replay arms.

The headline result is that phase 6 does not support the simple explanation that broader replay
coverage alone accounts for the phase-5 TSP outcome. The best mean arm is diversity-aware failure
replay, not random replay, and the archive diagnostics show
that hardness tracks final held-out TSPLIB gap more strongly than descriptor diversity does.

Table~\ref{tab:routing-phase6-summary} gives the most important mechanism comparisons.

\begin{table}[t]
\centering
\scriptsize
\setlength{\tabcolsep}{2pt}
\begin{tabularx}{\textwidth}{
  >{\raggedright\arraybackslash}p{1.45in}
  >{\raggedright\arraybackslash}p{0.85in}
  >{\raggedright\arraybackslash}p{1.08in}
  Y
}
\toprule
Comparison & Transfer delta & 95\% interval & Reading \\
\midrule
random replay vs. no replay & $-0.003123$ & $[-0.027964,0.021061]$ & slight mean advantage for random replay, but no decisive separation \\
raw failure replay vs. random replay & $+0.001361$ & $[-0.022536,0.025457]$ & naive raw-failure replay does not recover the phase-5 TSP win \\
diversity-aware failure replay vs. raw failure replay & $-0.008637$ & $[-0.035972,0.017862]$ & curation improves the mean arm, but the gain is not decisive at the paired level \\
residual-failure replay vs. raw failure replay & $+0.022578$ & $[-0.005643,0.04867]$ & residual-failure replay underperforms the naive raw-failure arm \\
random replay with compression vs. random replay & $+0.014874$ & not favorable & compression lowers novelty but worsens transfer on the random-replay arm \\
diversity-aware failure replay with compression vs. diversity-aware failure replay & $+0.021526$ & not favorable & compression also hurts the best diversity-aware replay arm \\
\bottomrule
\end{tabularx}
\caption{Key phase-6 TSP replay-mechanism comparisons. Positive deltas are worse because the metric
is lower-is-better transfer gap.}
\label{tab:routing-phase6-summary}
\end{table}

The diagnostic side of phase 6 sharpens the interpretation. Archive diversity has near-zero
association with held-out TSPLIB gap (Pearson $r=-0.009478$, Spearman $\rho=-0.041013$), while
archive hardness has a materially stronger positive association with worse held-out gap (Pearson
$r=0.39687$, Spearman $\rho=0.371779$). The mechanism story is therefore not that random replay
won because it maximized diversity in any simple sense. A more defensible reading is that naive
failure replay concentrated too heavily on hard cases, while curated replay could recover some of
that loss without producing a decisive universal win.

\section{Implications of the Routing Phases}
\label{sec:routing-implications}

Phases 5 and 6 refine the dissertation's replay claim substantially.

\begin{enumerate}[leftmargin=2em]
\item Replay can help on benchmark routing, but the effect is conditional rather than universal.
\item Raw failure replay is not a reliable default. Curation matters more than failure accumulation
alone.
\item Compression reduces novelty growth, but in the completed TSP mechanism study it does not
improve the winning replay arms.
\item The same replay logic does not transfer cleanly across TSP, ATSP, and CVRP.
\end{enumerate}

These findings are scientifically useful even though they are mixed. They prevent the thesis from
over-claiming, and they motivate the later phases that change the search object rather than simply
adding more replay variants to the same scaffold.
